\subsection{Synthetic Division: Method} 
\begin{Example}Use synthetic division to divide:
\[
\left(2x^3+x^2-13x-1\right)\div\left(x+3\right)
\]
\end{Example}
\begin{enumerate}
\item The divisor must be in the form $x-k$. $x+3=x-\makeblank[.5in]$, so $k=\makeblank[.5in]$.
\item In order to avoid subtraction, we \makeblank of the number that appears in the divisor. We write \makeblank[.5in] in the upper left corner and put a bracket around it.
\item Note: the number we write in synthetic division is equal to \makeblank[.5in]
\item Now, write the coefficients of the \makeblank, in order, to the right of that number.
\item Make a line underneath that row of numbers, with room for another row of numbers between.
\item Copy the \makeblank below the line. In our example, we copy the number \makeblank[.5in].
\item\label{synthstartloop} Multiply the number you just wrote by the number in the bracket. Write the result above the line in the next column.
\item\label{synthendloop} Add the numbers in that column, and write the result below the line.
\item Repeat steps \ref{synthstartloop} and \ref{synthendloop} until you have filled all the columns.
\item Make a vertical line before the last number.
\end{enumerate}
\begin{lpic}{}{3}
\end{lpic}
After we do synthetic division, we have to read the answer.
\begin{enumerate}
\item The degree of the quotient is one less than the degree of the \makeblank[1.25in]
\item In our example, the degree of the quotient is $\makeblank[.5in]-1=\makeblank[.5in]$.
\item So we start with \makeblank[.5in] and fill in the variables in the terms in descending order. The last number before the vertical line is the constant.
\item The number after the vertical line is the \makeblank.
\end{enumerate}
We have found that
\[
\left(2x^3+x^2-13x-1\right)\div\left(x+3\right)=\makeblank[3in]
\]